\documentclass[12pt]{article}
\usepackage{amsmath, amssymb}
\usepackage[margin=1in]{geometry}
\setlength{\parskip}{6pt}
\title{QUBO Formulation: Cryptanalysis Prime Factorization Toy Problem}
\date{}
\begin{document}
\maketitle

\subsection*{Cryptanalysis Prime Factorization Toy Problem}

\paragraph{Problem Statement.}
Given a composite integer $N$ and a finite set of $M$ candidate factor pairs $\{(p_i, q_i)\}_{i=0}^{M-1}$, the objective is to select exactly one pair whose product equals $N$. Each candidate is evaluated by a precomputed error metric $e_i = (p_i q_i - N)^2$, potentially augmented with additional penalties for trivial factors. The task is to choose the candidate that minimizes this error while satisfying the selection constraint.

\paragraph{Decision Variables.}
For each candidate $i \in \{0, 1, \dots, M-1\}$, introduce a binary variable $x_i \in \{0,1\}$. The variable $x_i = 1$ indicates that the $i$-th candidate factor pair is selected, and $x_i = 0$ otherwise.

\paragraph{QUBO Derivation.}
\textbf{Step 1 --- Objective.} The original optimization problem is to minimize the total error over the selected candidate:
\begin{align}
\min_{\mathbf{x} \in \{0,1\}^M} \sum_{i=0}^{M-1} e_i x_i.
\end{align}

\textbf{Step 2 --- Constraint.} The selection must satisfy the cardinality constraint that exactly one candidate is chosen:
\begin{align}
\sum_{i=0}^{M-1} x_i = 1.
\end{align}

\textbf{Step 3 --- Penalty Term Expansion.} We convert the equality constraint into a quadratic penalty term with weight $\lambda > 0$:
\begin{align}
P(\mathbf{x}) &= \lambda \left( \sum_{i=0}^{M-1} x_i - 1 \right)^2.
\end{align}
Expanding the square yields:
\begin{align}
P(\mathbf{x}) &= \lambda \left( \sum_{i=0}^{M-1} x_i^2 - 2\sum_{i=0}^{M-1} x_i + 1 \right).
\end{align}
Applying the idempotent property $x_i^2 = x_i$ for binary variables gives:
\begin{align}
P(\mathbf{x}) &= \lambda \left( \sum_{i=0}^{M-1} x_i - 2\sum_{i=0}^{M-1} x_i + 1 \right) = \lambda \left( 1 - \sum_{i=0}^{M-1} x_i \right) = \lambda - \lambda \sum_{i=0}^{M-1} x_i.
\end{align}
Thus, the penalty contributes a linear term $-\lambda x_i$ to each variable and a constant offset $\lambda$.

\textbf{Step 4 --- Combined Hamiltonian.} Adding the objective and the penalty term produces the full QUBO Hamiltonian:
\begin{align}
H(\mathbf{x}) &= \sum_{i=0}^{M-1} e_i x_i + \lambda - \lambda \sum_{i=0}^{M-1} x_i = \sum_{i=0}^{M-1} (e_i - \lambda) x_i + \lambda.
\end{align}

\textbf{Step 5 --- Q Matrix and Offset.} Reading off the coefficients from the general closed form of $H(\mathbf{x})$, the QUBO parameters are:
\begin{align}
Q_{i,i} &= e_i - \lambda \quad \text{for } i \in \{0, \dots, M-1\}, \\
Q_{i,j} &= 0 \quad \text{for } i \neq j, \\
c &= \lambda.
\end{align}

\paragraph{Penalty Weight Justification.}
Let $E_{\max} = \max_{i} e_i$. To guarantee that any feasible solution (exactly one $x_i=1$) has strictly lower energy than any infeasible solution (zero or multiple $x_i=1$), we require $\lambda > E_{\max}$. If zero variables are selected, the penalty evaluates to $\lambda$. If exactly one variable $x_i$ is selected, the energy is $e_i$. Since $e_i \le E_{\max} < \lambda$, all feasible solutions are energetically favored over the zero-selection case. If two or more variables are selected, the penalty term grows quadratically, yielding an energy of at least $\lambda$, which again exceeds any feasible energy. Therefore, $\lambda > E_{\max}$ ensures that constraint violations are strictly penalized.

\paragraph{Correctness Argument.}
Minimizing $H(\mathbf{x})$ over $\{0,1\}^M$ recovers the optimal factor pair because the penalty term enforces the cardinality constraint. Any feasible assignment has energy equal to its objective value, which is bounded above by $E_{\max} < \lambda$. Any infeasible assignment incurs a penalty of at least $\lambda$, making its energy strictly greater than that of any feasible assignment. Consequently, the global minimum of $H$ must correspond to a feasible solution, and among those, it selects the one with the smallest error $e_i$.

\paragraph{Illustrative Example.}
Consider an instance with $M=3$ candidates and precomputed errors $e_0 = 4$, $e_1 = 9$, and $e_2 = 1$. We choose $\lambda = 10$, which satisfies $\lambda > \max\{4, 9, 1\}$. Substituting these values into the general formulas yields the diagonal Q matrix entries $Q_{0,0} = 4 - 10 = -6$, $Q_{1,1} = 9 - 10 = -1$, and $Q_{2,2} = 1 - 10 = -9$, with $Q_{i,j} = 0$ for $i \neq j$. The offset is $c = 10$. The concrete Hamiltonian is:
\begin{align}
H(x_0, x_1, x_2) = -6x_0 - x_1 - 9x_2 + 10.
\end{align}
Enumerating the feasible assignments:
\begin{align}
H(1,0,0) &= 4, \quad H(0,1,0) = 9, \quad H(0,0,1) = 1.
\end{align}
Infeasible assignments (e.g., $x=(1,1,0)$ or $x=(0,0,0)$) yield energies $\ge 10$. The global minimum occurs at $x_2 = 1$ with energy $1$, correctly identifying the second candidate as the optimal factor pair.
\end{document}